\documentclass[12pt,a4paper]{article}
\usepackage[UTF8]{ctex}
\usepackage{amsmath}
\usepackage{amsfonts}
\usepackage{amssymb}
\usepackage{graphicx}
\usepackage{booktabs}
\usepackage{array}
\usepackage{geometry}
\geometry{left=2.5cm,right=2.5cm,top=2.5cm,bottom=2.5cm}
\usepackage{float}
\usepackage{longtable}

\title{薄透镜焦距的测定实验报告}
\author{学生签名：刘航宇}
\date{日期：2025年3月6日}

\begin{document}

\maketitle

\section{实验目的}
\begin{enumerate}
\item 深入理解自准法、物距像距法、共轭法测定凸透镜焦距的原理，熟练掌握其操作方法。
\item 精准掌握利用物距像距法测定凹透镜焦距的流程，深刻理解凹透镜成像特性及与凸透镜的差异。
\item 熟练且精准地调节光学系统共轴，显著提升实验数据记录、处理及分析误差的能力。
\end{enumerate}

\section{实验原理}

\subsection{凸透镜焦距测定}
\begin{enumerate}
\item \textbf{自准法}
\begin{itemize}
\item 原理：物屏位于凸透镜焦平面时，反射光经透镜成像于物屏，焦距 $f = |A_0 - O|$（$A_0$ 为物屏位置，$O$ 为透镜光心位置）。
\end{itemize}

\item \textbf{物距像距法}
\begin{itemize}
\item 原理：依据薄透镜成像公式 $\frac{1}{u} + \frac{1}{v} = \frac{1}{f}$，推导得 $f = \frac{uv}{u + v}$，通过测量物距 $u$、像距 $v$ 计算焦距。
\end{itemize}

\item \textbf{共轭法}
\begin{itemize}
\item 原理：物屏与像屏间距 $L > 4f$，透镜两次成像位置差为 $d$，则 $f = \frac{L^2 - d^2}{4L}$，减少直接测量物距像距的误差。
\end{itemize}
\end{enumerate}

\subsection{凹透镜焦距测定}
\begin{itemize}
\item \textbf{物距像距法}：借助凸透镜成像，凹透镜物距 $u$ 为虚物距，像距 $v$ 为虚像距（取负），通过 $f = \frac{uv}{u + v}$ 计算，遵循符号规则。
\end{itemize}

\section{实验仪器}

\begin{table}[H]
\centering
\begin{tabular}{|c|c|c|p{6cm}|}
\hline
仪器名称 & 规格/型号 & 数量 & 用途说明 \\
\hline
光具座 & 1.5m导轨 & 1 & 承载光学元件，提供测量基准 \\
\hline
凸透镜 & 焦距约10 cm & 1 & 用于凸透镜焦距测量及辅助凹透镜测量 \\
\hline
凹透镜 & 焦距约-10 cm & 1 & 测定凹透镜焦距 \\
\hline
平面反射镜 & 50mm×50mm & 1 & 配合自准法反射光线 \\
\hline
物屏（带箭孔） & 标准尺寸 & 1 & 提供清晰物像 \\
\hline
像屏（毛玻璃） & 标准尺寸 & 1 & 接收并显示像 \\
\hline
钠光灯 & 波长589.3nm & 1 & 提供稳定光源 \\
\hline
米尺 & 精度1mm & 1 & 测量距离 \\
\hline
\end{tabular}
\end{table}

\section{实验步骤}

\subsection{共轴调节}
\begin{enumerate}
\item \textbf{粗调}：安装元件，调节高度使中心大致同一高度。
\item \textbf{细调}：移动凸透镜，依据"大像追小像"原则，使像中心与物屏中心重合。
\end{enumerate}

\subsection{自准法测凸透镜焦距}
\begin{enumerate}
\item 固定物屏 $A_0 = 10.00 \, \text{cm}$，放置凸透镜与平面镜。
\item 移动凸透镜，记录清晰成像时透镜位置，左右读数各一次，重复测量。
\end{enumerate}

\subsection{物距像距法测凸透镜焦距}
\begin{enumerate}
\item 固定物屏，移动凸透镜，记录放大像、缩小像时透镜与像屏位置，计算 $u$、$v$ 及 $f$。
\end{enumerate}

\subsection{共轭法测凸透镜焦距}
\begin{enumerate}
\item 固定物屏与像屏间距 $L > 4f$，记录透镜两次成像位置，计算 $d$ 及 $f$。
\end{enumerate}

\subsection{物距像距法测凹透镜焦距}
\begin{enumerate}
\item 凸透镜先成像，记录像屏位置 $A_1$。
\item 插入凹透镜，重新记录像屏与凹透镜位置，计算 $u$、$v$ 及 $f$。
\end{enumerate}

\section{实验数据记录与处理}

\subsection{自准法测凸透镜焦距}

\begin{table}[H]
\centering
\begin{tabular}{|c|c|c|c|c|c|}
\hline
次数 & 凸透镜光心位置（左→右，cm） & 凸透镜光心位置（右→左，cm） & 平均光心位置（cm） & $f = |A_0 - O|$（cm） & $\overline{f}$（cm） \\
\hline
1 & 32.90 & 32.60 & 32.75 & 22.75 & 22.70 \\
\hline
2 & 32.70 & 32.60 & 32.65 & 22.65 & \\
\hline
\end{tabular}
\end{table}

\textbf{计算}：$\overline{f} = \frac{22.75 + 22.65}{2} = 22.70 \, \text{cm}$

\subsection{物距像距法测凸透镜焦距}

\begin{table}[H]
\centering
\begin{tabular}{|c|c|c|c|c|c|c|c|c|}
\hline
实像类型 & 凸透镜光心位置 $O$（cm） & 像屏位置 $A_1$（左→右，cm） & 像屏位置 $A_1$（右→左，cm） & 平均像屏位置（cm） & $u$（cm） & $v$（cm） & $f$（cm） & $\overline{f}$（cm） \\
\hline
放大像 & 40.00 & 132.30 & 132.20 & 132.25 & 30.00 & 92.25 & 22.64 & 22.67 \\
\hline
缩小像 & 80.00 & 113.30 & 113.90 & 113.60 & 70.00 & 33.60 & 22.70 & \\
\hline
\end{tabular}
\end{table}

\textbf{计算}：$\overline{f} = \frac{22.64 + 22.70}{2} = 22.67 \, \text{cm}$

\subsection{共轭法测凸透镜焦距}

\begin{table}[H]
\centering
\begin{tabular}{|c|c|c|c|c|c|}
\hline
次数 & 透镜位置 $O_1$（cm） & 透镜位置 $O_2$（cm） & $d$（cm） & $L$（cm） & $f$（cm） \\
\hline
1 & 88.55 & 41.30 & 47.25 & 110.00 & 22.52 \\
\hline
\end{tabular}
\end{table}

\textbf{计算}：$f = \frac{110^2 - 47.25^2}{4 \times 110} = 22.52 \, \text{cm}$

\subsection{物距像距法测凹透镜焦距}

\begin{table}[H]
\centering
\begin{tabular}{|c|c|c|c|c|c|c|}
\hline
次数 & 像屏位置 $A_1$（cm） & 凹透镜光心位置 $O_2$（cm） & $u$（cm） & $v$（cm） & $f$（cm） & $\overline{f}$（cm） \\
\hline
1 & 120.50 & 116.75 & 3.75 & -8.85 & -6.62 & -6.62 \\
\hline
\end{tabular}
\end{table}

\textbf{计算}：$f = \frac{3.75 \times (-8.85)}{3.75 - 8.85} = -6.62 \, \text{cm}$

\section{实验结果与分析}
\begin{enumerate}
\item \textbf{自准法}：$\overline{f} = 22.70 \, \text{cm}$，误差源于成像清晰度的主观判断。
\item \textbf{物距像距法}：$\overline{f} = 22.67 \, \text{cm}$，物距像距测量误差影响结果。
\item \textbf{共轭法}：$f = 22.52 \, \text{cm}$，通过间距测量减小误差，结果更精确。
\item \textbf{凹透镜}：$\overline{f} = -6.62 \, \text{cm}$，负号符合特性，误差来自共轴调节与像屏定位。
\end{enumerate}

\section{实验结论}
通过多种方法完成薄透镜焦距测定，掌握测量原理与方法。共轭法测凸透镜精度高，凹透镜测量需精细调节。误差主要来自仪器、操作与成像判断，可通过优化操作、多次测量减小误差。

\vspace{1cm}
\noindent 学生签名：刘航宇 \hfill 教师签名：

\vspace{0.5cm}
\noindent 日期：2025年3月6日

\end{document}